\subsection*{الف}
یک زنجیرهٔ پردازشی کامل (از صحنه تا نمایش) می‌تواند به‌صورت زیر نمایش داده شود:
\[
\boxed{\text{صفحهٔ چاپی}} \rightarrow
\boxed{\text{اپتیک}} \rightarrow
\boxed{\text{پیش‌پردازش آنالوگ}} \rightarrow
\boxed{\text{نمونه‌بردار}} \rightarrow
\boxed{\text{کوانتایزر}} \rightarrow
\boxed{\text{کُدگذاری دیجیتال}} \rightarrow
\boxed{\text{ذخیره}} \rightarrow
\boxed{\text{بافر فریم }} \rightarrow
\boxed{\text{نمایشگر}}
\]
در اینجا، چون خروجی نهایی دیجیتال نمایش داده می‌شود، مرحلهٔ \lr{D/A} معمولاً وجود ندارد و نمایشگر از بافر فریم دیجیتال تغذیه می‌شود.

\subsection*{ب}
برای تصویربرداری از صفحهٔ چاپی با خطوط نازک، فرمول‌ها و متن ریز، مناسب‌ترین روش دریافت، \lr{scan-in} است.

دلیل:
\begin{itemize}
    \item \lr{scan-in} معمولاً با یک آرایهٔ خطی و حرکت مکانیکیِ کنترل‌شده صفحه/هد انجام می‌شود؛ بنابراین می‌توان به نمونه‌برداری فضایی بسیار ریز و یکنواخت (به‌ویژه در راستای اسکن) رسید.
    \item هندسهٔ تصویر و فوکوس قابل‌کنترل‌تر است (صفحه تخت و فاصله ثابت)، و اعوجاج پرسپکتیو و تغییرات بزرگ‌نمایی مثل حالت دوربین کمتر رخ می‌دهد؛ این برای متن/فرمول که به لبه‌های تیز حساس است مهم است.
\end{itemize}

\subsection*{ج}
چرا دو روش دیگر مناسب‌تر نیستند:
\begin{itemize}
    \item \lr{self-scanning\ array}: در این روش (مثل \lr{CCD/CMOS} ناحیه‌ای) رزولوشن محدود به تعداد پیکسل‌های آرایه و کیفیت اپتیک است؛ برای متن ریز، یا باید سنسور بسیار پرپیکسل/اپتیک دقیق داشت یا ریسکِ محوشدگی، آستیگماتیسم/اعوجاج لنز، و \lr{moir\'e} روی الگوهای چاپی بالا می‌رود.
    \item \lr{scan-out}: این روش (به‌معنای اسکن‌کردن تصویر روی سنسور/خروجی به‌صورت رستری، در سامانه‌های قدیمی) معمولاً با محدودیت‌های ساختاری مثل غیرخطی‌بودن اسکن، اعوجاج هندسی/هم‌خطی، و محدودیت پهنای‌باند روبه‌روست؛ در نتیجه برای نمایش دقیق لبه‌های تیز و جزئیات ریزِ سند انتخاب خوبی نیست.
\end{itemize}

\subsection*{د}
اسکن درهم‌بافته (\lr{interlaced}) انتخاب خوبی نیست.

دو دلیل فنی:
\begin{itemize}
    \item برای محتوای سند (خطوط نازک و متن)، نمونه‌برداری عمودیِ درهم‌بافته باعث \lr{line twitter} / سوسو و آلیاسینگ روی جزئیات افقی می‌شود و جزئیات ریز بین دو \lr{field} ناپایدار دیده می‌شود.
    \item نمایشگرهای دیجیتال امروز ذاتاً \lr{progressive} هستند؛ بنابراین باید \lr{deinterlacing} انجام شود که معمولاً یا لبه‌ها را نرم/محو می‌کند یا مصنوعات زمانی-فضایی تولید می‌کند؛ هر دو برای متن و فرمول نامطلوب‌اند.
\end{itemize}

\subsection*{هـ}
مقایسهٔ \lr{NTSC}، \lr{PAL/SECAM} و \lr{CIF}:
\begin{itemize}
    \item \lr{NTSC}: حدود $525$ خط با $60$ \lr{field/s} و ساختار درهم‌بافته؛ رزولوشن مؤثر عمودی حدود $480$ خطِ فعال. به‌دلیل درهم‌بافتگی و محدودیت‌های ویدئوی آنالوگ (پهنای‌باند محدود، مصنوعات ترکیب رنگ)، برای متن/فرمول ضعیف است.
    \item \lr{PAL/SECAM}: حدود $625$ خط با $50$ \lr{field/s} و درهم‌بافته؛ رزولوشن عمودیِ فعال معمولاً حدود $576$ خط، پس از نظر تعداد خطوط از \lr{NTSC} بهتر است؛ اما همچنان مشکل اصلی، درهم‌بافتگی و ماهیت ویدئویی/آنالوگ بودن و نمایش نامناسب لبه‌های تیز است.
    \item \lr{CIF}: معمولاً $352 \times 288$ پیکسل (دیجیتال). مزیتِ مهم: دیجیتال بودن و (در بسیاری پیاده‌سازی‌ها) فریم‌های \lr{progressive} و پیکسل‌های نسبتاً منظم؛ اما رزولوشن بسیار پایین است و برای متن ریز/فرمول‌ها معمولاً ناکافی می‌شود.
\end{itemize}

جمع‌بندی: هیچ‌کدام واقعاً ایده‌آلِ تصویربرداری سند نیستند. اگر مجبور به انتخاب بین این سه مرجع تاریخی باشیم:
\begin{itemize}
    \item از نظر رزولوشن فضایی خام، \lr{PAL/SECAM} معمولاً از \lr{NTSC} معقول‌تر است (خطوط فعال بیشتر).
    \item از نظر ساختار اسکن و سازگاری با نمایش دیجیتال (بدون درهم‌بافتگی/ترکیب آنالوگ)، \lr{CIF} از نظر «ساختار» بهتر است، ولی به‌خاطر $288$ خط عمودی و عرض کم، برای متن/فرمول مناسب نیست.
\end{itemize}
پس در عمل برای این کاربرد باید سراغ قالب‌های اسکن سند با رزولوشن بالاتر و \lr{progressive} رفت، نه استانداردهای ویدئویی تاریخی.